\documentclass[aspectratio=169,10pt,english]{beamer}

\input{../../../_preambulo_beamer.tex}
\input{../../../_comandos_beamer.tex}

\selectlanguage{english}

\title{MA1001B - Statistical Modeling for Decision Making}
\subtitle{Lesson 22: Normal Applications: Tolerances and Operating Intervals}
\author{}
\institute{Tecnol\'ogico de Monterrey}
\date{}

\begin{document}

\begin{frame}[plain]
  \vspace{1.2cm}
  {\Large\bfseries MA1001B - Statistical Modeling for Decision Making\par}
  \vspace{0.7cm}
  {\large Lesson 22: Normal Applications and Tolerances\par}
  \vspace{0.7cm}
  {\color{TecRojo}\rule{\textwidth}{1.2pt}}
  \vspace{0.8cm}

  MA1001B Faculty\\[0.3cm]
  Term\\[0.3cm]
  Tecnol\'ogico de Monterrey
\end{frame}

\begin{frame}{The Specification Question}
  \begin{alertblock}{Guiding question}
    If fill weight is $N(50,4)$ grams and the spec is $[44,56]$, what
    fraction of units is in spec, and what happens if the mean moves?
  \end{alertblock}

  \vspace{0.6em}
  By the end of this lesson, you can:
  \begin{itemize}
    \item convert spec limits into $z$ cuts;
    \item compute yield as $P(\mathrm{LSL}<X<\mathrm{USL})$;
    \item quantify the yield drop after a mean shift;
    \item explain why high yield is not the same as being centered.
  \end{itemize}

  \vfill
  \scriptsize All fill weights used today are synthetic. No real company data are used.
\end{frame}

\begin{frame}{Centered Process and Spec Cuts}
  \small
  Model $X\sim N(50,4^2)$. Specification $[44,56]$:
  \[
    z_{\mathrm{LSL}}=\frac{44-50}{4}=-1.5,\qquad
    z_{\mathrm{USL}}=\frac{56-50}{4}=1.5.
  \]
  In-spec yield:
  \[
    P(44<X<56)=0.866386.
  \]
  Each tail is $0.066807$ because the process is centered.

  \vspace{0.3em}
  \begin{exampleblock}{Synthetic filling line}
    The spec is a customer interval. Yield is the area between the cuts.
  \end{exampleblock}
\end{frame}

\begin{frame}[fragile]{Step 1: In-Spec Yield}
  \begin{columns}[T]
    \begin{column}{0.42\textwidth}
      \small
      \begin{lstlisting}
z_low = (44 - 50) / 4
z_high = (56 - 50) / 4
yield_p = (
    model.cdf(56)
    - model.cdf(44)
)
      \end{lstlisting}

      \begin{block}{Verified output}
        $z=\pm 1.500000$\\
        Each tail $=0.066807$\\
        Yield $=0.866386$
      \end{block}
    \end{column}
    \begin{column}{0.58\textwidth}
      \centering
      \includegraphics[width=\textwidth]{../figures/normal_tolerances_01_spec_yield.png}
      \vspace{0.2em}
      \scriptsize Centered yield from Step 1
    \end{column}
  \end{columns}
\end{frame}

\begin{frame}{A Mean Shift Moves the $z$ Cuts}
  \small
  The spec stays at $[44,56]$. After $\mu=52$,
  \[
    z_{\mathrm{LSL}}=\frac{44-52}{4}=-2,\qquad
    z_{\mathrm{USL}}=\frac{56-52}{4}=1.
  \]
  Yield falls to $0.818595$. The drop versus the centered process is
  $0.047791$. The upper tail grows to $0.158655$.
\end{frame}

\begin{frame}[fragile]{Step 2: Yield After $\mu=52$}
  \begin{columns}[T]
    \begin{column}{0.40\textwidth}
      \small
      \begin{lstlisting}
shifted = norm(52, 4)
new_yield = (
    shifted.cdf(56)
    - shifted.cdf(44)
)
      \end{lstlisting}

      \begin{block}{Verified output}
        $z$ cuts $-2$ and $1$\\
        Yield $=0.818595$\\
        Drop $=0.047791$
      \end{block}
    \end{column}
    \begin{column}{0.60\textwidth}
      \centering
      \includegraphics[width=\textwidth]{../figures/normal_tolerances_02_mean_shift.png}
      \vspace{0.2em}
      \scriptsize Centered versus shifted densities from Step 2
    \end{column}
  \end{columns}
\end{frame}

\begin{frame}{Step 3: High Yield Need Not Be Centered}
  \begin{columns}[T]
    \begin{column}{0.42\textwidth}
      \small
      Tight off-center clock: $\mu=52$, $\sigma=2$.

      \vspace{0.4em}
      $z$ cuts $-4$ and $2$.\\
      Yield $=0.977218$, higher than the centered $0.866386$.

      \vspace{0.4em}
      \textbf{Limit:} spec compliance is not the same as being centered.
      Lesson 23 studies the sampling distribution of $\bar x$.
    \end{column}
    \begin{column}{0.58\textwidth}
      \centering
      \includegraphics[width=\textwidth]{../figures/normal_tolerances_03_centered_vs_compliant.png}
      \vspace{0.2em}
      \scriptsize Centered-wide versus tight off-center from Step 3
    \end{column}
  \end{columns}
\end{frame}

\begin{frame}{Concept Check}
  \begin{enumerate}
    \item Compute the $z$ cuts for spec $[44,56]$ on $N(50,4)$.
    \item Why are the two tails equal to $0.066807$ in Step 1?
    \item After $\mu=52$, which tail grows, and to what value?
    \item Why can yield $0.977218$ still be the wrong operating story?
  \end{enumerate}

  \vspace{0.6em}
  \begin{block}{Connection to Lesson 23}
    The session can continue directly into the sampling distribution of
    the sample mean and its standard error.
  \end{block}

  \vfill
  \scriptsize Source: OpenStax, \textit{Introductory Business Statistics 2e},
  Chapter 6, Section 6.2. CC BY-NC-SA.
\end{frame}

\appendix



\end{document}
